\section{Preliminaries} \label{02_modeling}
This section describes the problem we studied and preliminary
knowledge of world models.
%\subsection{Problem Definition}
In real-world driving, an ego vehicle cannot access the complete and exact state of the environment, such as the precise intentions of other drivers or occluded road regions. Instead, it must rely solely on its onboard sensors, which provide partial and noisy observations. We model the autonomous driving task as a Partially Observable Markov Decision Process (POMDP).  The POMDP is defined by the tuple $(\mathcal{S}, \mathcal{A}, \mathcal{O}, \mathcal{T}, \mathcal{G}, \mathcal{R}, \gamma)$, where  $\mathcal{S}$ represents the true but unobservable environment state. $\mathcal{A}$ is the action space consisting of continuous steering and throttle/brake commands $\mathbf{a}_t = (a_{\text{steer}}, a_{\text{throttle}}) \in [-1, 1]^2$. $\mathcal{O}$ is the observation space. At each time step $t$, the ego vehicle receives an observation $\mathbf{o}_t$ composed of two modalities:
\begin{equation}
\mathbf{o}_t = \{I_t, \mathbf{v}_t\},
\end{equation}
where $I_t \in \mathbb{R}^{H \times W \times 3}$ is an image from a front-facing camera, and $\mathbf{v}_t \in \mathbb{R}^5$ is a vector of vehicle physics, including speed, steering angle, previous actions and yaw rate. Both of them can be directly obtained through raw sensors or physical calculations. $\mathcal{T}(s' \mid s, a): \mathcal{S} \times \mathcal{A} \times \mathcal{S} \rightarrow [0, 1]$ defines the conditional transition probability distribution over next states. $\mathcal{G}(o \mid s, a): \mathcal{S} \times \mathcal{A} \times \mathcal{O} \rightarrow [0, 1]$ defines the probability of observing given state and action. $\mathcal{R}: \mathcal{S} \times \mathcal{A} \rightarrow \mathbb{R}$ is the reward function balancing progress, safety, and rule compliance, while $\gamma \in [0, 1]$ is the discount factor. The agent's objective is to learn a policy $\pi(a_t | o_{\leq t})$ that maximizes the expected cumulative discounted reward:
\begin{equation}
\mathbb{E}_{\pi} \left[ \sum_{t=0}^{\infty} \gamma^t r_t \right].
\end{equation}
%\subsection{World Models}
Due to partial observability in driving, the true transition dynamics $\mathcal{T}$ are unknown. To address this challenge, world-model-based reinforcement learning approximates belief dynamics using a learned latent representation. Instead of explicitly maintaining a probability distribution over environment states, the agent learns a compact latent variable $\hat{s}_{t} \in \mathcal{S}_{\text{latent}}$ that summarizes past information.

The world model usually includes the following components: a representation model $p(\hat{s}_{t} \mid \hat{s}_{t-1}, a_{t-1}, o_{t})$ that maps high-dimensional observations into a compact latent state; a latent transition model $p(\hat{s}_{t+1} \mid \hat{s}_t, a_t)$ that predicts future latent states; an observation model that reconstructs or predicts observations $p(\hat{o}_{t} \mid \hat{s}_{t})$; and a reward model $p(\hat{r}_{t} \mid \hat{s}_{t}, a_t)$ for value estimation. In some tasks, there also has a continue model $p(\hat{c}_{t} \mid \hat{s}_{t}, a_t) $ that predicts termination signal. Through this learned latent dynamics model, the agent can perform imagination rollouts entirely in latent space, enabling policy optimization without requiring real environment interaction at every step.